
%----------------------------------------------------------------------------------------
%	DOCUMENT DEFINITIONS
%----------------------------------------------------------------------------------------

\documentclass[a4paper]{resume}

\usepackage[left=0.75in,top=0.6in,right=0.75in,bottom=0.6in]{geometry}
\usepackage[T1]{fontenc}
\usepackage[utf8]{inputenc}	
\usepackage[brazil]{babel}
\usepackage{microtype}
\usepackage{lmodern}
\usepackage{enumitem}
\usepackage{multicol}
\usepackage{booktabs}
\usepackage{hyperref}

\hypersetup{
    colorlinks=true,
    linkcolor=blue,
    filecolor=magenta,      
    urlcolor=black,
    pdftitle={Pedro Henrique Rosso - Currículo},
    }

\urlstyle{same}

\newcommand{\tab}[1]{\hspace{.2667\textwidth}\rlap{#1}}
\newcommand{\itab}[1]{\hspace{0em}\rlap{#1}}

\name{Pedro Henrique Di Francia Rosso}
\address{Rua Almirante Barroso, 1095 (Apto 1104), Centro, Criciúma-SC, Brasil, CEP 88802-249}
\address{+55 (48) 99956-8129 \\ {\tt pedrohrosso@gmail.com | p233687@dac.unicamp.br}}

\begin{document} 

%----------------------------------------------------------------------------------------
%	EDUCATION SECTION
%----------------------------------------------------------------------------------------

\begin{rSection}{Formação}
{\bf Bacharelado em Engenharia de Computação} \hfill {\em 2014--2018} \\
{\sc Universidade Federal de Santa Catarina (UFSC)} \hfill {\em Araranguá-SC, Brasil} \\
Participei de maratonas de programa, fui monitor de disciplinas e bolsista de iniciação 
científica~\cite{pibic}. Finalizei a gradução estudando recoloração de imagens para dicromatas~\cite{tcc}. 
Também fui o melhor aluno da turma e entre engenharias (Computação e Energia) com coeficiente 9.04.

{\bf Mestrado em Ciências da Computação} \hfill {\em 2019--2021} \\
{\sc Universidade Federal do ABC (UFABC)} \hfill {\em Santo André-SP, Brasil} \\
Completei o curso com coeficiente A em todas disciplinas cursadas. Defendi a
dissertação intitulada ``OCFTL: {\it an MPI implementation-independent
fault tolerance library for task-based
applications}'' em julho de 2021 sob orientação do Prof. Emilio Francesquini.

{\bf Doutorado em Ciência da Computação} \hfill {\em 2021--Atual} \\
{\sc Universidade Estadual de Campinas (Unicamp)} \hfill {\em Campinas-SP, Brasil} \\
Atualmente estou cursando Doutorado em Ciência da Computação na Universidade
Estadual de Campinas, com enfoque em Tolerância a Falhas para OpenMP/MPI. 
Sob orientação do Prof. Guido Araújo e Prof. Emilio Francesquini (UFABC).
\end{rSection}

%----------------------------------------------------------------------------------------
%	SKILLS SECTION
%----------------------------------------------------------------------------------------

\begin{rSection}{Prêmios e Distinções}

{\bf Menção Honrosa} $\diamond$ {UFSC} \hfill{\em 2015--2018} \\
Durante os 4 últimos semestres da graduação, recebi a Menção Honrosa pela UFSC por 
obter coefiente médio de notas do semestre acima de 9.0. Do 6$^o$ ao 9$^o$ semestre.

{\bf Melhor Desempenho Acadêmico} $\diamond$ {UFSC} \hfill{\em 2018} \\
Em feveireiro de 2019, recebi uma medalha da Universidade Federal de Santa Catarina por obter o
melhor desempenho acadêmico no curso de Bacharelado em Engenharia de Computação -- Turma 2018.

{\bf Melhor Aluno entre Engenharias} $\diamond$ {CREA-SC} \hfill{\em 2018} \\
Em feveireiro de 2019, recebi uma distinção do Conselho Regional de Engenharia e Agronomia
por obter o melhor desempenho acadêmico entre as Engenheiras de Computação e Energia -- Turma 2018.

{\bf Prêmios em Artigos} $\diamond$ {ERAD-SP} \hfill{\em 2020 e 2021} \\
Nas ERAD-SP dos anos de 2020 e 2021, obtive o prêmio de melhor artigo na trilha de Pós-Gradução.

{\bf Menção Honrosa} $\diamond$ {WSCAD-CTD} \hfill{\em 2021} \\
No WSCAD-CTD de 2021, minha dissertação de mestrado recebeu menção honrosa por estar entre as três finalistas
do concurso.

\end{rSection}

\begin{rSection}{Interesses}
A seguir uma lista não exaustiva de interesses profissionais e de pesquisa:

\setlength\multicolsep{5pt}
\begin{multicols}{2}
\begin{itemize}[noitemsep]
  \item Computação de Alto Desempenho;
  \item Tolerância a Falhas;
  \item Compiladores;
  \item Computação Paralela e Distribuída;
  \item Arquitetura de Computadores.
\end{itemize}
\end{multicols}

\end{rSection}

%----------------------------------------------------------------------------------------
%	EXPERIENCE SECTION
%----------------------------------------------------------------------------------------

\begin{rSection}{Experiência}

{\bf Monitoria de Estrutura de Dados} \hfill{\em 2016} \\
{\sc Universidade Federal de Santa Catarina (UFSC)} \hfill {\em Araranguá-SC, Brasil} \\
Durante o 1$^o$ e 2$^o$ semestres de 2016 realizei acompanhamento dos alunos da disciplina 
de Estruturas de Dados, onde realizei revisões de conteúdo, exercícios e provas. Monitoria aprovada com nota máxima.

{\bf Bolsista de Iniciação Científica} \hfill{\em 2017--2018} \\
{\sc Universidade Federal de Santa Catarina (UFSC)} \hfill {\em Araranguá-SC, Brasil} \\
Durante o 2$^o$ semestre de 2017 e 1$^o$ semestre de 2018 participei de um projeto de Iniciação Científica, onde foi desenvolvida uma aplicação de eletromiografia sem fio. Iniciação Científica financiada pelo CNPq.

{\bf Bolsista de Mestrado} \hfill{\em 2019--2021} \\
{\sc Universidade Federal do ABC (UFABC)} \hfill {\em Santo André-SP, Brasil} \\
Do ano de 2019 até 2021 realizei minha pesquisa de mestrado com enfoque em tolerância a falhas para MPI. Trabalho de mestrado financiado pela FUNCAMP. Durante o mesmo periodo, também fui participante do projeto OmpCluster sob coordenação do prof. Guido Araújo (UNICAMP).

{\bf Bolsista de Doutorado} \hfill{\em 2021--Atual} \\
{\sc Universidade Estadual de Campinas (Unicamp)} \hfill {\em Campinas-SP, Brasil} \\
Desde o ano de 2021 realizo de doutorado com enfoque em tolerância a falhas para OpenMP/MPI. Trabalho de doutorado financiado pela FAPESP (2021/09355-2). Participante projeto OmpCluster sob coordenação do prof. Guido Araújo (UNICAMP).

\end{rSection}


%----------------------------------------------------------------------------------------
%	SKILLS SECTION
%----------------------------------------------------------------------------------------

\begin{rSection}{Habilidades e Proficiência em Línguas}
{\bf Linguagens de Programação, Frameworks}

\begin{tabular}{ll}
  Maior Experiência & C, C++ \\
  Experiência Média & Java, JavaScript, Node.js\\
  Possuo familiaridade & Assembly, C\#, Bash, CUDA, Javascript, Robótica \\
\end{tabular}

{\bf Ferramentas}
\begin{itemize}[noitemsep]
  \item Sistema operacional Linux e programas utilitários ({\tt grep}, {\tt
    bash}, {\tt make}, etc);
  \item Paradigmas de programação paralela OpenMP e MPI;
  \item Controle de versionamento ({\tt git}) e plataformas (Github, Gitlab,
    etc);
  \item Sistemas embarcados (projeto, prototipagem e software).
  \item Programação em MatLab/Octave.
  \item Debuggers ({\tt gdb}) e perfiladores ({\tt perf}, {\tt strace}, {\tt
    ltrace});
  \item Infraestrutura de compiladores LLVM.
  \item Compiladores {\tt gcc} e {\tt clang} de C/C++;
\end{itemize}

{\bf Proficiência em Línguas}
\begin{itemize}[noitemsep]
  \item {\bf Português} (Nativo) -- compreendo bem; falo bem; escrevo bem.
  \item {\bf Inglês} (Avançado) -- compreendo bem; falo razoável; escrevo bem.
\end{itemize}
\end{rSection}

\begin{rSection}{Participação em Eventos}


{\bf 12$^o$ CARLA} $\diamond$ México, Virtual \hfill {\em 2021} \\
Autor~\cite{carla2021}. Latin America High Performance Computing Conference.

{\bf 12$^a$ ERAD-SP} $\diamond$ UFABC/USP, Virtual \hfill {\em 2021} \\
Autor~\cite{erad2021}. Escola Regional de Alto Desempenho do Estado de São Paulo.

{\bf 11$^a$ ERAD-SP} $\diamond$ UNESP/Mackenzie/USP, Virtual \hfill {\em 2020} \\
Autor~\cite{erad2020}. Escola Regional de Alto Desempenho do Estado de São Paulo.

{\bf 28$^o$ SIC} $\diamond$ UFSC, Florianópolis-SC, Brazil \hfill {\em 2018} \\
Autor~\cite{pibic}. Seminário de Iniciação Científica.

{\bf 7$^o$ SICT-SUL} $\diamond$ UFSC, Araranguá-SC, Brazil \hfill {\em 2018} \\
Autor~\cite{pibic}. Simpósio de Integração Científica e Tecnológica do Sul Catarinense.

\end{rSection}
\newpage
\begin{rSection}{Links}
  \begin{tabular}{ll}
    Homepage       & \url{https://pedroohr.github.io} \\
    Github         & \url{https://github.com/pedroohr} \\
    Gitlab         & \url{https://gitlab.com/phrosso} \\
    LinkedIn       & \url{https://twitter.com/pedrohrosso} \\
    Google Scholar & \url{https://scholar.google.com/citations?user=rtONezgAAAAJ} \\
    Lattes         & \url{http://lattes.cnpq.br/5343894554617060} \\
  \end{tabular}
\end{rSection}

\begin{rSection}{Publicações}
\renewcommand{\section}[2]{}
\bibliographystyle{abbrv}
\bibliography{reference}
\end{rSection}

\vfill \hfill {\em Atualizado em: \today.}

\end{document}